\documentclass[12pt, a4paper]{article}
\usepackage[spanish]{babel}             % Español para los nombres
\usepackage[utf8]{inputenc}             % Codificación de entrada
\usepackage[T1]{fontenc}                % Codificación de fuente (escribir castellano)
\usepackage{lmodern}                    % Fuente (la default no es compatible con castellano)
\usepackage{csquotes}                   % Para comillas y otras anotaciones del castellano
\usepackage{fancyhdr}                   % Encabezado y Pie de pagina
\usepackage{tocloft}                    % Puntos en los índices

\usepackage{hyperref}                   % Links y referencias dentro del texto
\usepackage{acronym}                    % Para incluir las listas de abreviaturas
\usepackage{multicol}                   % Permite crear espacios con columnas verticales
\usepackage{parskip}                    % Eliminar sangrado y añadir espacio entre párrafos

\usepackage{graphicx}                   % Para incluir imágenes
\usepackage{subcaption}                 % Subfiguras
\usepackage{xcolor}                     % Definir y utilizar colores
\usepackage{tikz}                       % Dibujar formas, figuras, rectas, intersecciones...
\usepackage[absolute,overlay]{textpos}  % Posición absoluta para textos
\usepackage{geometry}                   % Para controlar los márgenes

\usepackage{float}


\usepackage[backend=biber, sorting=none]{biblatex} % Referencias (bibliografía / webgrafía)
\selectlanguage{spanish}                % Seleccionar español
\bibliography{referencias}              % Incluir el archivo de referencias

% Definición de nombres
%\printbibliography[title={Bibliografía}]
\renewcommand{\cftsecleader}{\cftdotfill{\cftdotsep}} 
\renewcommand{\spanishtablename}{Tabla}
\renewcommand{\spanishlisttablename}{Índice de tablas}
\newcommand{\estudiante}{García, Alan}
\newcommand{\director}{xxxxx xx, xxxx}
\newcommand{\grado}{INGENIERÍA XXX}
\newcommand{\titulo}{PLANTILLA TFG PARA LA ESCUELA DE INGENIERÍA DE BILBAO}
\newcommand{\portada}{images/protada.png}
\newcommand{\curso}{2023-2024}

% Definición de Macros
\newcommand{\aclink}[1]{\hyperlink{acro:#1}{\ac{#1}}} % Enlazar la abreviatura al acrónimo

% Encabezado y Pie de pagina
\pagestyle{fancy}                       % Estilo de las páginas
\fancyhf{}
\fancyhead[L]{Trabajo Fin de Grado}
\fancyhead[R]{\rightmark}
\fancyfoot[L]{UPV/EHU}
\fancyfoot[R]{\thepage}
\renewcommand{\footrulewidth}{0.4pt} % Línea horizontal en el pie de página
\setlength{\headheight}{15.5pt}

% Definición de colores
\definecolor{ehu_blue}{HTML}{376092}
\definecolor{link_color}{HTML}{36AEB4}
\definecolor{reference_color}{HTML}{0F3133}

% Configuración de hyperref para diferentes tipos de enlaces
\hypersetup{
    colorlinks=true,
    linkcolor=reference_color,
    citecolor=reference_color,
    filecolor=link_color,
    urlcolor=link_color,
}

% #########################################################################
% #                                  TFG                                  #
% #########################################################################
\begin{document}
\include{title} % Incluimos el título
\newpage
\input{resumen} % Incluimos el resumen tri-lingüe

% Índices
\tableofcontents\thispagestyle{empty}\newpage
\listoffigures\thispagestyle{empty}\newpage
\listoftables\thispagestyle{empty}\newpage

% Abreviaturas
\addcontentsline{toc}{section}{Abreviaturas}
\section*{Abreviaturas}
\input{abreviaturas}

\section{Introducción}
En 1983 el autor James Watson descubrió el \aclink{ADN} \cite{watson53}. La <<medicina>> realmente era veneno. El \aclink{ASSI} es el encargado de mostrar en todo momento el estado del \aclink{AS}. Si el \aclink{ASSI} se apaga de forma inesperada, se considerará un estado de emergencia eléctrico.
El \aclink{USB} es un protocolo de comunicación basado en el protocolo RS232.

En la wikipedia de LaTex \cite{latex_wikibook} puedes encontrar información muy útil para tus documentos.

\subsection{Estudio demográfico}
En los últimos años, la población mundial ha ido decayendo considerablemente. Tanto es así, que los países más desarrollados, los cuales son a su vez los más envejecidos, están teniendo muchas dificultades para mantener sus sistemas de pensiones. Estos hechos se ven reforzados por los datos del gráfico \ref{fig:envejecimiento} en el que se comparan distintas poblaciones Europeas \footnote{Aguacate con salami}.

\begin{figure}[H]
    \centering
    \includegraphics[width=5cm, height=3cm]{images/ejemplo.png}
    \caption{Envejecimiento en Europa}
    \label{fig:envejecimiento}
\end{figure}


\begin{table}[htbp]
    \centering
    \begin{tabular}{|c|c|c|}
        \hline
        \textbf{Nombre} & \textbf{Edad} & \textbf{Género} \\
        \hline
        Juan & 25 & Masculino \\
        María & 30 & Femenino \\
        Carlos & 22 & Masculino \\
        Ana & 28 & Femenino \\
        \hline
    \end{tabular}
    \caption{Ejemplo de tabla}
    \label{tab:ejemplo}
\end{table}
\newpage
¿Qué pasa si sigo escribiendo aquí? En principio debería poner en el encabezado la sección en la que estamos, pero no estoy seguro. Al parecer, funciona perfecto.

\newpage\printbibliography
\addcontentsline{toc}{section}{Referencias}
\end{document}
